\documentclass[book.tex]{subfiles}

\begin{document}

\chapter{Differentiable Convex Optimization}
\label{chap:deep_equilibrium}
\noindent\rule{11cm}{0.4pt} \\
\textbf{Prerequisites:}  \autoref{chap:ml_basics}, \autoref{chap:probability}  \\
\textbf{Difficulty Level:} ** \\
\noindent\rule{11cm}{0.4pt}
\vspace{5mm}

Convex optimization problems provide a convenient mathematical language for stating a class of mathematical problems that are amenable to optimization methods. These layers can now be used as layers in differentiable programs, facilitating modeling of physical systems with constrained degrees of freedom.

\section{Convex Optimization and Convex Programs}
A general family of convex optimization problems can be written as

\begin{align}
    x^*(\theta) = &\arg\min \quad f_0(x; \theta) \\
    &\textrm{subject to} \quad f_i(x; \theta) \leq 0,\quad i = 1,\dotsc,m \\
    &\qquad \qquad \quad  g_j(x;\theta) = 0, \quad j=1,\dotsc, n
\end{align}

Here $x$ is the optimization variable, $f_0, f_i, g_j$ are all convex functions and $\theta$ is a parameter vector.

Convex programs have increasingly become standard computational primitives that can be used as pieces of more complex programs. Can we use a convex problem as a part of a differentiable program? It is in fact possible to differentiate through convex optimization problems \cite{agrawal2019differentiable}, \cite{agrawal2020learning}. If the objective and condition functions are smooth, we can write down the KKT equations that govern the solution and solve them directly 

At times, we will have non-smooth, non-differentiable objective or condition functions. In this case, we can implicitly define the stationary conditions of a generic cone program.

\begin{align}
\textrm{minimize} \quad c^T x \quad
\textrm{subject to} \quad b - A x \in \mathcal{K}
\end{align}
where $\mathcal{K}$ is a convex cone, $A$ is a matrix, and $b$ is a vector.

%\begin{lstlisting}[language=python]
%x = Variable()
%objective = (x - 2 $\theta$)$^2$
%problem = Minimize(objective)
%\end{lstlisting}

\section{Disciplined Convex Programming}
The field of disciplined convex programming has introduced domain specific languages to faciliate specification of complex convex programs. These systems, such as cvxpy, perform compilation steps to transform an arbitrary convex optimization into a standardized cone program. When computing gradients, these transformations also have to be differentiated through as well. The compilation transformations are affine so the differentiation is simple.

A differentiable disciplined convex program can be constructed by using differentiable disciplined components in a Physika program.

\section{Taking Derivatives of a Cone Program}

Suppose that we have a cone program given in primal form by
\begin{align}
    &\textrm{minimize}\ c^T x \\
    & \textrm{subject to}\ Ax + s = b \\
    & \qquad \qquad \ s \in \mathcal{K}
\end{align}

We swap to dual formulation of cone program. View a conic program as a function

\begin{align}
    \psi: \mathbb{R}^{m \times n}\times \mathbb{R}^m \times \mathbb{R}^n \to \mathbb{R}^{n+2m}
\end{align}
We compute the adjoint to the derivative of $\psi$ at $(A, b, c)$ by the formula
\begin{align}
    (dA, db, dc) &= D^T \psi(A, b, c)(dx, dy, ds) \\
    &= D^T Q(A, b, c)D^T s(Q) D^T \phi(z)(dx, dy, ds)
\end{align}


\section{Noisy Linear Regression}

Suppose $(x_i, y_i)$ are the data. Suppose we have a linear classifier

\begin{align}
    \hat{y} &= 1[\beta^T x \geq 0]
\end{align}
The loss is given by
\begin{align}
    \mathcal{L} &= \frac{1}{m} \sum \ell(\theta, x_i, y_i)
\end{align}
$\beta^*$ is the solution to the convex optimization problem for training. 

%\textcolor{red}{TODO: The original talk cites a data poisoning example. Perhaps this can be reframed as a broken instrument problem}

The gradient $\nabla_{x_i} \mathcal{L}$ gives the direction that creates a maximal increase in the training loss.




\end{document}